\documentclass{article}

\usepackage{amsmath,amssymb}
\usepackage{xurl}
\usepackage[colorlinks=true,linkcolor=blue,citecolor=blue,urlcolor=blue]{hyperref}
\setlength{\emergencystretch}{3em}

\input{./command}

\title{The CPSV16 Unit-Weight Convention and Renormalization-Fixed-Point Scale}
\author{}
\date{2026-08-15}

\begin{document}
\maketitle
\gapnote{false-source}{resolved}

\section{Two scale conventions in the source}

The canonical-form construction of
Cirac--P\'erez-Garc\'ia--Schuch--Verstraete first writes a tensor as
\begin{align}
    A^i
    &=\bigoplus_k \mu_k A_k^i,
    \label{eq:uw_canonical_decomposition}
\end{align}
where every $A_k$ is normal and its transfer map has spectral radius one
\cite[lines~214--246]{Cirac2016MPDO_arXiv}. Since the states are not normalized
there, line~246 adopts the standing convention
\begin{align}
    |\mu_k|
    &\leq 1
    &&\text{for every $k$},
    \label{eq:uw_weight_upper_bound}\\
    |\mu_{k_0}|
    &=1
    &&\text{for some $k_0$}.
    \label{eq:uw_unit_weight_witness}
\end{align}
The formal BNT predicate records
Eq.~(\ref{eq:uw_unit_weight_witness}) through
\leanid{MPSTensor.IsBNTCanonicalForm.weight_unit_exists}. This condition gives
one global unit-weight witness, not one witness in every sector. The latter
distinction is discussed separately in
\path{docs/paper-gaps/cpsv16_global_vs_persector_unit_witness.tex}.

Definition~4.1 of Ref.~\cite{Cirac2016MPDO_arXiv} fixes the tensor scale
differently. For a tensor $M$ already in canonical form and generating MPDOs,
it requires trace-preserving completely positive maps relating the one-site
and two-site physical closures
\cite[lines~645--660]{Cirac2016MPDO_arXiv}. Define the physical-trace transfer
by
\begin{align}
    \mathcal T_M
    &=\sum_i M^{ii}.
    \label{eq:uw_physical_trace_transfer}
\end{align}
Trace preservation gives
\begin{align}
    \mathcal T_M^2
    &=\mathcal T_M,
    \label{eq:uw_physical_idempotence}
\end{align}
as calculated at lines~1333--1340 of
Ref.~\cite{Cirac2016MPDO_arXiv}. The source-facing formal statement is
\leanid{MPOTensor.isPhysicalTraceIdempotent_of_isRFPViaTS}; its proof uses
\leanid{MPOTensor.physTraceTransfer_sq_of_isRFPViaTS}.

These requirements concern two transfer objects with different scaling laws.
Define the doubled-index MPS transfer map by
\begin{align}
    \mathcal E_M(X)
    &=\sum_i M^iX(M^i)^\dagger.
    \label{eq:uw_mps_transfer_definition}
\end{align}
For a scalar $c$,
\begin{align}
    \mathcal E_{cM}
    &=|c|^2\mathcal E_M.
    \label{eq:uw_mps_transfer_scaling}
\end{align}
Linearity of the physical-trace transfer gives
\begin{align}
    \mathcal T_{cM}
    &=c\mathcal T_M.
    \label{eq:uw_physical_transfer_scaling}
\end{align}
Equations~(\ref{eq:uw_weight_upper_bound}) and
(\ref{eq:uw_unit_weight_witness}) normalize the first object by choosing normal
BNT representatives and placing the overall scale in their coefficients.
Equation~(\ref{eq:uw_physical_idempotence}) constrains the second object at the
scale of the displayed tensor. The related up-to-scalar treatment of source
zero correlation length is recorded in
\path{docs/paper-gaps/cpsv16_zcl_canonical_form_normalization.tex}.

\section{Fixed-representative scale pinning}

By Eq.~(\ref{eq:uw_physical_transfer_scaling}), replacing $M$ by $cM$ replaces
$\mathcal T_M$ by $c\mathcal T_M$. Suppose that
$\mathcal T_M^2=\mathcal T_M\ne0$ and that the nonzero rescaling $cM$ satisfies
the same literal idempotence equation. Then
\begin{align}
    c^2\mathcal T_M
    &=(c\mathcal T_M)^2
    \label{eq:uw_rescaled_transfer_square}\\
    &=c\mathcal T_M,
    \label{eq:uw_rescaled_idempotence}
\end{align}
so $c=1$. Thus Definition~4.1 of
Ref.~\cite{Cirac2016MPDO_arXiv} is not invariant under arbitrary nonzero
rescaling of a fixed representative.

By contrast, the simplicity definition at lines~815--822 of
Ref.~\cite{Cirac2016MPDO_arXiv} is phrased through the elements of a BNT: the
tensor is simple when none of those elements has nilpotent ket-against-bra
contraction. Multiplying one BNT element by a nonzero scalar does not change
nilpotency. This observation concerns only the nilpotency clause; it does not
imply that arbitrary complex rescaling preserves the MPDO condition or either
full simplicity predicate.

Let $\rho_N(M)$ denote the closed length-$N$ MPO. Its exact scaling law is
\begin{align}
    \rho_N(cM)
    &=c^N\rho_N(M),
    \label{eq:uw_closed_mpo_scaling}
\end{align}
not $|c|^{2N}\rho_N(M)$. Consequently, strictly positive real rescaling $rM$
preserves positivity in both directions. The formal theorems
\leanid{MPOTensor.isMPDO_smul_ofReal_iff} and
\leanid{MPOTensor.isSimple_smul_ofReal_iff} establish the corresponding
invariances for the MPDO and simplicity predicates. They do not assert
invariance under arbitrary nonzero complex scalars, and they do not imply that
the rescaled tensor satisfies the same literal fixed-scale RFP equations of
Definition~4.1 of Ref.~\cite{Cirac2016MPDO_arXiv}.

The declaration \leanid{MPOTensor.IsSimple} formalizes lines~217--246 and
Definition~4.7 of Ref.~\cite{Cirac2016MPDO_arXiv}. It asks for a positive
blocking length and a BNT sector presentation whose representative
physical-trace transfers are all non-nilpotent. Every representative has a
positive number of copies, and every copy has nonzero weight, following the
nonzero-coefficient convention discussed in
\path{docs/paper-gaps/cpsv16_bnt_uniqueness_zero_coefficient.tex}. The theorem
\leanid{MPOTensor.IsSimple.exists_mpo_ne_zero} derives a positive length $N$
such that
\begin{align}
    \rho_N(M)
    &\ne0,
    \label{eq:uw_positive_length_nonzero}
\end{align}
from the presentation; this is not an additional defining assumption.
Individual lengths may still vanish.

The predicate does not impose the line-246 unit-weight convention. Requiring a
unit-weight witness for the blocked tensor would be a hyper-literal reading
excluded by the source's own conventions. The dimer $R$ considered below is
simple in the source's sense because its physical-trace transfer is
non-nilpotent, and it carries canonical-form weight
$\sqrt{337/512}$. It would therefore not be simple under a unit-weight reading.
Example~4.12 of Ref.~\cite{Cirac2016MPDO_arXiv} is not a second such witness:
its second BNT basis element has vanishing one-by-one physical-trace transfer,
so it is not simple under either reading, and its canonical-form weight
$1/\sqrt2$ is beside the point. See
\path{docs/paper-gaps/cpsv16_simple_tensor_nilpotency.tex}. The library
therefore treats line~246 as an available source normalization, not as a clause
of Definition~4.7. The predicates that once separated these readings, and the
counterexample module that separated them, were removed; see
\path{docs/audits/2026-08-24_degenerate_readings_wave_2.md}.

By contrast, \leanid{MPOTensor.IsSimpleCanonicalForm} is a predicate on one
displayed canonical-form representative and supplies the Appendix~C.2
hypothesis of Ref.~\cite{Cirac2016MPDO_arXiv} for an already-blocked tensor. It
is not a quotient by nonzero scalar rescaling, and no equivalence with
\leanid{MPOTensor.IsSimple} is asserted.

\section{The parity tensor in Example~4.12}

Example~4.12 of Ref.~\cite{Cirac2016MPDO_arXiv} prints the two physical letters
$M^{00}=I$ and $M^{11}=\sigma_z$. Its physical-trace transfer satisfies
\begin{align}
    \mathcal T_M^2
    &=2\mathcal T_M,
    \label{eq:uw_parity_transfer_quasi_idempotence}
\end{align}
so the two trace-preserving channel equations cannot hold for the printed
representative. Define the density-normalized representative by
\begin{align}
    \widehat M
    &=\tfrac12 M.
    \label{eq:uw_parity_density_normalization}
\end{align}
Its physical-trace transfer is idempotent. Explicit parity coarse-graining and
refinement channels prove
\leanid{MPOTensor.CPSVExample412NormalizedRFP.Mhat_isRFPViaTS} without any
additional hypothesis; the complete calculation appears in
\path{docs/paper-gaps/cpsv16_example_4_12_normalization.tex}.

The doubled-index MPS has two bond-one sectors, with physical vectors
$(1,0,0,1)$ and $(1,0,0,-1)$. Each vector has squared norm $2$, so division by
$\sqrt2$ gives its normal representative. Consequently, the two horizontal
normal-block weights of $\widehat M$ are both $1/\sqrt2$. The scale selected by
the channel equations therefore does not supply the global unit-weight witness
in Eq.~(\ref{eq:uw_unit_weight_witness}). The formal theorem deliberately
records the bare channel predicate and does not claim that $\widehat M$
satisfies every standing canonical-form convention in Definition~4.1 of
Ref.~\cite{Cirac2016MPDO_arXiv}.

\section{The displayed dimer representative}

The project tensor $R$ gives a concrete instance of this tension. Its
doubled-index MPS transfer map satisfies
\begin{align}
    \mathcal E_R^2
    &=\frac{337}{512}\mathcal E_R.
    \label{eq:uw_dimer_transfer_quasi_idempotence}
\end{align}
Consequently, its displayed one-block CPSV canonical form is
\begin{align}
    R^{\bullet\bullet}
    &=\mu\widehat R,
    \label{eq:uw_dimer_canonical_form}\\
    \mu
    &=\sqrt{\frac{337}{512}},
    \label{eq:uw_dimer_weight}
\end{align}
where the transfer map of the retained block $\widehat R$ is idempotent and
therefore has spectral radius one. This presentation has $0<\mu<1$, so it does
not supply the global unit-weight witness in
Eq.~(\ref{eq:uw_unit_weight_witness}).

The physical-trace transfer of the displayed tensor $R$ is a nonzero
idempotent; its $(0,0)$ entry is $1/2$. The retained block satisfies
\begin{align}
    \mathcal T_{\widehat R}
    &=\mu^{-1}\mathcal T_R.
    \label{eq:uw_dimer_retained_transfer}
\end{align}
Since $\mu\ne1$, the scale-pinning calculation in
Eqs.~(\ref{eq:uw_rescaled_transfer_square}) and
(\ref{eq:uw_rescaled_idempotence}) shows that this normalized representative
does not satisfy the same literal physical-trace idempotence equation as $R$.

The obstruction is stronger than this displayed-presentation mismatch. For
every positive blocking length $L$, the doubled-index transfer map of the
blocked tensor satisfies
\begin{align}
    \mathcal E_{R^{[L]}}^2
    &=\left(\frac{337}{512}\right)^L\mathcal E_{R^{[L]}}.
    \label{eq:uw_blocked_quasi_idempotence}
\end{align}
Suppose that $R^{[L]}$ admitted a normalized horizontal BNT witness. The gauge
to that witness transports Eq.~(\ref{eq:uw_blocked_quasi_idempotence}). The
global unit-weight condition from line~246 of
Ref.~\cite{Cirac2016MPDO_arXiv} supplies a copy whose weight has modulus one.
Restricting the transported equation to that copy removes the weight. Its
normal representative is left canonical, so its transfer map is trace
preserving. Taking the trace of the quasi-idempotence equation at the identity
forces
\begin{align}
    \left(\frac{337}{512}\right)^L
    &=1.
    \label{eq:uw_blocked_unit_contradiction}
\end{align}
This contradicts
\begin{align}
    0
    &<\left(\frac{337}{512}\right)^L
    <1.
    \label{eq:uw_blocked_strict_bound}
\end{align}
Hence no positive blocking of $R$ admits a normalized horizontal canonical
form: no rescaling of $R$ reconciles the line-246 normalization with the scale
pinned by Definition~4.1 of Ref.~\cite{Cirac2016MPDO_arXiv}.

This argument is not a nilpotency statement and does not bear on
Definition~4.7 of Ref.~\cite{Cirac2016MPDO_arXiv}. The retained-block theorem
proves that the displayed canonical presentation has non-nilpotent
physical-trace transfer. Together with positivity and the one-site BNT
refinement, it witnesses \leanid{MPOTensor.IsSimple}\,$R$. The theorem
\leanid{R_isSimple} in the namespace
\leanid{MPOTensor.RescalingStableLengthDependentRFP} proves this assertion.
The source quantifier is existential and is witnessed by $L=1$; no assertion
is made about every positive blocking. Reading line~246 into Definition~4.7
would instead make $R$ non-simple, which is the reading rejected here.

\section{Coefficient absorption with all Appendix~C.2 conditions}

The two-sector coefficient-absorption example realizes the normalization
tension without dropping the conditions used in the simple Case-II passage
\cite[lines~1628--1665 and 1740--1782]{Cirac2016MPDO_arXiv}. Its nonzero
physical slices are
\begin{align}
    K^{00}
    &=K^{11}
    =\operatorname{diag}(1/2,0),
    \label{eq:uw_caseii_equal_slices}\\
    K^{22}
    &=\operatorname{diag}(0,1).
    \label{eq:uw_caseii_third_slice}
\end{align}
For $N>0$,
\begin{align}
    \bra{\sigma}\rho_N(K)\ket{\tau}
    &=\delta_{\sigma,\tau}\left(2^{-N}\mathbf 1_{\{\sigma_k\in\{0,1\}\ \forall k\}}
      +\mathbf 1_{\{\sigma_k=2\ \forall k\}}\right),
    \label{eq:uw_caseii_mpo_entries}\\
    \tr\rho_N(K)
    &=2>0.
    \label{eq:uw_caseii_mpo_trace}
\end{align}
The tensor therefore generates positive semidefinite periodic operators with
nonzero normalization.

Define
\begin{align}
    P
    &=\operatorname{diag}(1,1,0),
    \label{eq:uw_caseii_projection_p}\\
    Q
    &=I-P,
    \label{eq:uw_caseii_projection_q}\\
    \omega_P
    &=\tfrac14(P\otimes P),
    \label{eq:uw_caseii_state_p}\\
    \omega_Q
    &=Q\otimes Q.
    \label{eq:uw_caseii_state_q}
\end{align}
The one-site and two-site physical closures are
\begin{align}
    M_1(X)
    &=\tfrac{X_{00}}2P+X_{11}Q,
    \label{eq:uw_caseii_one_site_closure}\\
    M_2(X)
    &=X_{00}\omega_P+X_{11}\omega_Q.
    \label{eq:uw_caseii_two_site_closure}
\end{align}
Define the partial-trace and measure-and-prepare maps by
\begin{align}
    \mathcal S(Y)
    &=\tr_2(Y),
    \label{eq:uw_caseii_coarse_graining_map}\\
    \mathcal R(Y)
    &=\tr(PYP)\omega_P+\tr(QYQ)\omega_Q.
    \label{eq:uw_caseii_refinement_map}
\end{align}
These maps are trace preserving and completely positive, and they satisfy
\begin{align}
    \mathcal S[M_2(X)]
    &=M_1(X),
    \label{eq:uw_caseii_coarse_graining_identity}\\
    \mathcal R[M_1(X)]
    &=M_2(X).
    \label{eq:uw_caseii_refinement_identity}
\end{align}
Equations~(\ref{eq:uw_caseii_coarse_graining_identity}) and
(\ref{eq:uw_caseii_refinement_identity}) are the explicit Definition~4.1
equations from lines~638--660 of Ref.~\cite{Cirac2016MPDO_arXiv}. They give
SAL, while the slice sum gives the literal Definition~4.2 identity
\begin{align}
    \mathcal T_K
    &=I_2,
    \label{eq:uw_caseii_transfer_identity}\\
    \mathcal T_K^2
    &=\mathcal T_K.
    \label{eq:uw_caseii_transfer_idempotence}
\end{align}

The displayed representatives are normal and biCF, the weights
$(1/\sqrt2,1)$ satisfy the global unit-weight convention, and the ambient
tensor has normalized fixed-representative simple canonical form.
Nevertheless,
\begin{align}
    \rho(\mathcal E_{\mu_0A_0})
    &=\frac12,
    \label{eq:uw_caseii_absorbed_spectral_radius}\\
    \mu_0A_0
    &\text{ is not normal}.
    \label{eq:uw_caseii_absorbed_not_normal}
\end{align}
The component results and their conjunction are in
\path{TNLean/MPS/MPDO/CaseIIAbsorptionCounterexample.lean}, culminating in
\leanid{MPOTensor.CaseIIAbsorptionCounterexample.%
ambient_isSAL_isSimpleCanonicalForm_and_firstAbsorbed_not_isNormalTensor}.
This repairs the previously missing full-condition counterexample statement
and classifies the printed absorption inference as \emph{proof-path drift}. It
does not decide the claimed all-sector factorization.

The non-Cartesian witness in
\path{TNLean/MPS/MPDO/NonCartesianActiveSectorCounterexample.lean} refutes the
implication from injectivity, SAL, literal ZCL, and a scalar relation to a
normal tensor. Let $\mu$ be its singleton coefficient and let $r$ be the
Perron value of its unnormalized transfer, as defined in
\path{docs/paper-gaps/cpgsv17_pf_rank_one.tex}. They satisfy
\begin{align}
    |\mu|^2
    &=r,
    \label{eq:uw_noncartesian_weight_perron_relation}\\
    0
    &<|\mu|
    =\sqrt r
    <1.
    \label{eq:uw_noncartesian_strict_subunit_weight}
\end{align}
The strict bound in Eq.~(\ref{eq:uw_noncartesian_strict_subunit_weight}) is
formalized by
\leanid{MPOTensor.NonCartesianActiveSectorCandidate.%
transferMap_posDef_eigenvalue_lt_one}. Thus the unit clause of the line-246
convention has no witness. Rescaling to unit weight destroys literal ZCL.
Accordingly, \ghissue{6775} remains open at the ambient source-context
boundary.

This open issue is also the remaining input to the source's channel
construction. Proposition~\texttt{3to5} of
Ref.~\cite{Cirac2016MPDO_arXiv} constructs the representative maps from the
neighboring-trace factorization, and Proposition~\texttt{prop2to5} combines
them after projecting onto the orthogonal outer subspaces. The latter
combination is formalized in \ghissue{6632}. The direct sector-mixing
alternative formerly tracked in \ghissue{6793} is not a separate source step.

\section{Unequal raw copies refute Theorem~4.9(iv)$\Rightarrow$(v)}

The scale mismatch is decisive even when the source's global unit-weight
convention is satisfied. Let $A=1$ be the sole scalar BNT representative, and
take two raw copies with weights $1$ and $1/2$. The one-letter tensor then has
local virtual matrix
\begin{align}
    K^{00}
    &=\operatorname{diag}(1,1/2).
    \label{eq:uw_repeated_copy_local_matrix}
\end{align}
This is a BNT canonical presentation: both weights have modulus at most one,
and the first has modulus one. Every positive-length contraction is
$1+2^{-N}>0$, the scalar representative is normal and non-nilpotent, and the
singleton family has simultaneous one-letter span. Condition~(iv) of
Theorem~4.9 in Ref.~\cite{Cirac2016MPDO_arXiv} also holds: the representative
generates MPDOs, the distinct-layer equation is vacuous, and the scalar
physical-sector factorization has a positive neighboring operator and
normalized rank-one trace coefficients.

The sole local matrix of $K^{[2]}$ is $(K^{00})^2$. If the blocked tensor
satisfied Definition~4.1 of Ref.~\cite{Cirac2016MPDO_arXiv}, its
physical-trace transfer would be idempotent, so
\begin{align}
    (K^{00})^4
    &=(K^{00})^2.
    \label{eq:uw_repeated_copy_required_idempotence}
\end{align}
Taking traces would give
\begin{align}
    1+2^{-4}
    &=1+2^{-2},
    \label{eq:uw_repeated_copy_false_trace_equality}
\end{align}
which is false. The declaration in
\path{TNLean/MPS/MPDO/Theorem49RepeatedCopyCounterexample.lean},
\leanid{MPOTensor.CaseIIAbsorptionCounterexample.%
printed_theorem49_iv_to_v_is_false}, kernel-checks explicit component witnesses
for all standing data, MPDO positivity of the blocked tensor, and the failure
in Eq.~(\ref{eq:uw_repeated_copy_false_trace_equality}) in one conjunction.
Thus the literal implication (iv)$\Rightarrow$(v) in Theorem~4.9 of
Ref.~\cite{Cirac2016MPDO_arXiv} is refuted and classified as an
\emph{unfaithful statement}. The printed route exhibits \emph{proof-path
drift}: condition~(iv) does not constrain repeated-copy weights, so no
outer-sector assembly can repair the claim.

The viable source route is instead (ii)$\Rightarrow$(v). In the earlier
Case-II argument, lines~1646--1665 of Ref.~\cite{Cirac2016MPDO_arXiv} use
literal zero correlation length to make all copy weights within each BNT
sector equal and absorb their common value. The repeated-copy tensor in
Eq.~(\ref{eq:uw_repeated_copy_local_matrix}) fails condition~(ii), because
$\operatorname{diag}(1,1/2)$ is not idempotent. Consequently, it does not
contradict Proposition~\texttt{prop2to5} of
Ref.~\cite{Cirac2016MPDO_arXiv}. The remaining work is either the printed
per-representative factorization or another construction of the same blocked
channel pairs, followed by their projector-controlled combination after
common-weight absorption.

\section{The vertical one-label identity}

The same scale bookkeeping appears in the algebraic form of Theorem~4.14 of
Ref.~\cite{Cirac2016MPDO_arXiv}. Its length-one condition is
\begin{align}
    m_\gamma
    &=\sum_{\alpha,\beta}
      c^{(1)}_{\alpha,\beta,\gamma}m_\alpha m_\beta,
    \label{eq:uw_vertical_length_one_identity}
\end{align}
where
\begin{align}
    m_\alpha
    &=\tr(\mu_\alpha).
    \label{eq:uw_vertical_weight_trace}
\end{align}
These formulas appear at lines~972--993 of
Ref.~\cite{Cirac2016MPDO_arXiv}. In the displayed one-label dimer
presentation,
\begin{align}
    c^{(1)}_{0,0,0}
    &=\frac{32}{25},
    \label{eq:uw_dimer_structure_coefficient}\\
    m_0
    &=\frac{25}{32}.
    \label{eq:uw_dimer_vertical_weight}
\end{align}
Thus the nonzero solution is fixed by $(c^{(1)}_{0,0,0})^{-1}$ rather than by a
unit-weight convention. This statement concerns the explicit presentation and
is not an invariance theorem for coefficients under changes of presentation.

\section{Horizontal periodic equality does not fix the vertical normalization}

The distinction between a horizontal periodic family and a chosen vertical
presentation is already visible for one physical letter. Define
\begin{align}
    P
    &=\begin{pmatrix}1&0\\0&0\end{pmatrix},
    \label{eq:uw_horizontal_projection_p}\\
    X
    &=\begin{pmatrix}1&-3/4\\0&1\end{pmatrix},
    \label{eq:uw_horizontal_gauge_x}\\
    Q
    &=XPX^{-1}
    =\begin{pmatrix}1&3/4\\0&0\end{pmatrix}.
    \label{eq:uw_horizontal_projection_q}
\end{align}
Let $M_P$ and $M_Q$ be the one-letter MPO tensors with sole matrices $P$ and
$Q$, respectively. Since $P$ and $Q$ are idempotent and have trace one, every
positive length $N$ satisfies
\begin{align}
    \rho_N(M_P)
    &=\rho_N(M_Q)
    =1.
    \label{eq:uw_same_horizontal_family}
\end{align}
Thus both tensors generate MPDOs and have the same unnormalized periodic
operator family. They are related by the nonunitary horizontal gauge $X$,
which belongs to the general MPV gauge freedom described at lines~187--199 of
Ref.~\cite{Cirac2016MPDO_arXiv}.

Their normalized one-label vertical representatives differ:
\begin{align}
    A_P
    &=P,
    \label{eq:uw_vertical_representative_p}\\
    m_P
    &=1,
    \label{eq:uw_vertical_normalization_p}\\
    A_Q
    &=\frac45Q,
    \label{eq:uw_vertical_representative_q}\\
    m_Q
    &=\frac54.
    \label{eq:uw_vertical_normalization_q}
\end{align}
Both are normal bond-one tensors, and each gives a positive one-label algebra
clause. Their products satisfy
\begin{align}
    A_P^2
    &=A_P,
    \label{eq:uw_vertical_product_p}\\
    A_Q^2
    &=\frac45A_Q.
    \label{eq:uw_vertical_product_q}
\end{align}
Consequently, their structure coefficients are
\begin{align}
    c_P^{(L)}
    &=1,
    \label{eq:uw_vertical_coefficients_p}\\
    c_Q^{(L)}
    &=\left(\frac45\right)^L.
    \label{eq:uw_vertical_coefficients_q}
\end{align}
A singleton label set has only one relabelling, so these coefficients cannot
agree after relabelling. The first family is length independent, whereas the
second is not.

The declarations in
\path{TNLean/MPS/MPDO/VerticalCoefficientPresentationCounterexample.lean}
prove the gauge relation, MPDO positivity, both tensor-attached algebra
clauses, Eqs.~(\ref{eq:uw_vertical_coefficients_p}) and
(\ref{eq:uw_vertical_coefficients_q}), and the failure of relabelled equality.
In particular,
\leanid{MPOTensor.VerticalCoefficientPresentationCounterexample.%
sameMPV₂Pos_and_no_relabelled_coefficient_equality}
refutes the theorem proposed in \ghissue{6395} under its stated bare
positive-length MPV-equality hypothesis.

This example is not a counterexample to Proposition~4.13 or Theorem~4.14 of
Ref.~\cite{Cirac2016MPDO_arXiv}. The source starts with one horizontally
canonical MPDO tensor and then chooses its vertical decomposition
\cite[lines~943--991]{Cirac2016MPDO_arXiv}. No horizontal canonical-form
assertion is made here for the sheared tensor $M_Q$. The source comparison in
Appendix~C.4 retains the positive ratio
$m_\alpha/n_{\sigma(\alpha)}$ instead of discarding the normalization scale
\cite[lines~1997--2008]{Cirac2016MPDO_arXiv}. Consequently, the source's
fixed-presentation characterization and the project's explicit
length-dependent RFP example remain unchanged. This counterexample does not
establish rigidity between two literal horizontal canonical presentations.

A scale-covariant statement requires explicit vertical transport. Suppose
there are a label equivalence $e$, nonzero scalars $s_\alpha$, and algebra
isomorphisms $\Phi_L$ such that every label $\alpha$ satisfies
\begin{align}
    O'_L(\alpha)
    &=s_\alpha^L\Phi_L(O_L(e\alpha)).
    \label{eq:uw_vertical_explicit_transport}
\end{align}
At a positive length where the target operators are linearly independent,
uniqueness of the product expansion gives
\begin{align}
    {c'}^{(L)}_{\alpha,\beta,\gamma}
    &=\left(\frac{s_\alpha s_\beta}{s_\gamma}\right)^L
      c^{(L)}_{e\alpha,e\beta,e\gamma}.
    \label{eq:uw_vertical_coefficient_covariance}
\end{align}
Exact equality therefore requires a normalization condition such as
$s_\alpha s_\beta=s_\gamma$ on every active product triple. Deriving
Eq.~(\ref{eq:uw_vertical_explicit_transport}) from two horizontal canonical
presentations, and extending
Eq.~(\ref{eq:uw_vertical_coefficient_covariance}) to every positive length,
are separate statements not supplied by bare horizontal periodic equality.

\bibliographystyle{alpha}
\bibliography{references}

\end{document}
